\chapter{Conclusions}
\label{chap:conclusions}

\section{Summary}
\label{sec:conclusion-summary}

\subsubsection*{The Initial Problem}

This thesis addressed the gap between GPU allocation and actual GPU use in Kubernetes-based AI inference. Accelerators are expensive and energy-intensive resources, yet a workload that reserves an entire device may use only a fraction of its capacity. The initial objective was therefore to increase GPU utilization through sharing. During the work, however, this objective evolved into a more meaningful question: \textbf{how much useful application work can be completed with the available hardware, time, and energy?} High utilization alone is not evidence of efficiency if it does not translate into more completed requests, shorter completion times, or lower energy per useful output.

The investigation combined \textbf{NVIDIA MIG}, \textbf{Kubernetes DRA}, and workload-level orchestration. MIG was used to divide the accelerator into isolated resources, while DRA provided a declarative way to request them according to their properties. Two proofs of concept then examined different allocation lifecycles: a long-running Ollama inference service and a pool of finite Jobs managed through Kueue.

\subsubsection*{What Was Implemented and Observed}

In the Ollama experiment, restricting a single server to a smaller MIG partition reduced its individual performance. Right-sizing is therefore not free, even for a lightweight model. When several isolated replicas were deployed and sufficient request concurrency was available, however, the same physical GPU executed more useful work in parallel. The Multiple Instances configuration increased aggregate \textbf{User Throughput by 114.05\%} relative to the single-server baseline under high load and reduced the energy required per generated token. The benefit came primarily from \textbf{MIG partitioning, service replication, and workload consolidation}; DRA improved provisioning and removed the need to bind applications to specific device identifiers, but each claim was resolved only when its long-running Pod started.

The Kueue proof of concept focused instead on finite workloads, whose resources are requested and released at every execution. GPU-aware service classes combined admission priority with DRA eligibility constraints, allowing each newly admitted Job to use either a medium or a large compatible partition. Compared with a capacity-equivalent static mixed policy, D-Flex reduced the pool \textbf{makespan by 14.59\%}. DRA did not make an individual GPU computation faster; its contribution was to keep multiple compatible resources selectable until allocation time, thereby reducing unused capacity at the end of the queue.

Together, the experiments show two distinct advantages. MIG can turn unused capacity into parallel, isolated execution when the workload tolerates smaller partitions. DRA can make the allocation policy more flexible when resource requests recur and cannot be optimally fixed in advance. \textbf{The benefit is produced by the complete system design---not by MIG or DRA in isolation.}

\subsubsection*{Main Considerations}

The main lessons of the study can be summarized as follows:

\begin{itemize}
    \item \textbf{Useful work matters more than utilization.} GPU activity is a diagnostic signal, not an objective by itself. It must be interpreted together with User Throughput, latency, makespan, and energy per unit of work. The time-sharing experiment, for example, kept the accelerator active without providing a comparable throughput improvement.

    \item \textbf{Energy efficiency must be evaluated per useful output.} An idle GPU retains a non-negligible base power draw, so consolidating compatible workloads can avoid leaving several accelerators powered but lightly used. At the same time, a configuration may draw more instantaneous power and still consume less total energy when it completes substantially more work in a shorter interval.

    \item \textbf{The appropriate sharing policy depends on the workload.} AI inference includes persistent language-model services, embedding and anomaly-detection pipelines, and finite classification or batch tasks. Their memory demand, concurrency, duration, and isolation requirements differ. The choice among time-sharing, MPS, MIG, and dedicated devices must therefore follow the workload rather than a universal infrastructure rule.

    \item \textbf{Allocation lifecycle determines the value of DRA.} A long-running service generally acquires a device once and then serves many application requests without creating new resource claims. Finite Jobs repeatedly release and request devices, giving DRA new opportunities to select among compatible resources.

    \item \textbf{``Dynamic'' does not mean autonomous.} DRA does not inspect an application, predict its demand, or independently choose the ideal GPU partition. Administrators still define DeviceClasses, selectors, quotas, and service-class policies. DRA supplies the mechanism to express and enforce those decisions consistently at allocation time.
\end{itemize}

The resulting infrastructure should hide physical identifiers and placement details from applications without ignoring meaningful hardware differences. \textbf{The central conclusion is not that GPU sharing should always be increased, but that each platform should maximize useful work while balancing performance, energy efficiency, isolation, flexibility, and operational complexity for the workload lifecycle at hand.}

\section{Future Work}
\label{sec:future-work}

The experiments in this thesis used a MIG geometry configured before each test. The administrator, through MIG Manager, retained ownership of the partition layout, while DRA selected among the instances already present. Recent NVIDIA DRA driver developments introduce a stronger interpretation of dynamic allocation through the alpha \texttt{DynamicMIG} feature. Instead of requiring every partition to be created in advance, the driver advertises the valid MIG profiles and placements of a partitionable GPU. When a claim selects a compatible placement and sufficient capacity is available, the node-side driver creates the requested GPU and compute instances during resource preparation. It destroys them again when the claim is released~\cite{nvidia2026drivermig}.

This mechanism shifts part of the partition-lifecycle ownership from the cluster administrator to the workload request, mediated by Kubernetes and the NVIDIA driver. The administrator still defines the allowed DeviceClasses and operating policy, but no longer needs to predict one fixed geometry for all subsequent workloads. Such a model could reduce stranded capacity when the required mix of partition profiles changes over time and could make service classes portable across nodes with different current layouts.

Dynamic MIG should not be interpreted as transparent live elasticity. It does not resize or migrate the device of a running CUDA process, and a new allocation decision still requires an appropriate workload or Pod lifecycle boundary. Physical placement constraints can also prevent a requested profile from being created even when the aggregate amount of free compute or memory appears sufficient. Since the feature is currently alpha and disabled by default, evaluating its provisioning latency, cleanup behaviour, fragmentation, recovery after failures, and interaction with admission control represents a direct continuation of this work. On architectures where MIG mode cannot be enabled without resetting the GPU, that mode must also remain an administrator-managed prerequisite.

A longer-term development would add a workload-aware decision layer above DRA. Such a component could estimate resource demand from the model, numerical precision, input dimensions, batch size, expected duration, concurrency, and service-level objective. Telemetry from previous executions could refine the estimate and determine whether the next workload should request a small isolated partition, a larger profile, or a different sharing mechanism. The prediction would then be translated into a service class or DRA selector before allocation, while Kueue would continue to enforce admission policy and quota.

Combining this prediction layer with Dynamic MIG would approach the final objective of automatically selecting and materializing the most appropriate hardware resource at runtime. The decision would initially remain scoped to safe lifecycle boundaries, such as Job submission, retry, or replica creation, rather than attempting to resize active processes in place. A feedback loop could subsequently compare predicted and observed completion time, throughput, memory demand, and energy per useful output, improving later allocations. This evolution would preserve the abstraction introduced by DRA while moving the policy from manually defined static expectations toward adaptive, evidence-based resource selection.
